\begin{figure*}[t]
\centering
\begin{tikzpicture}[
    node distance=0.4cm,
    box/.style={
        rectangle,
        draw=tscgblue!80,
        fill=tscgblue!8,
        rounded corners=2pt,
        minimum height=0.7cm,
        minimum width=1.8cm,
        font=\scriptsize\sffamily,
        align=center,
        thick
    },
    iobox/.style={
        rectangle,
        draw=tscggray!80,
        fill=tscggray!10,
        rounded corners=2pt,
        minimum height=0.7cm,
        minimum width=1.8cm,
        font=\scriptsize\sffamily,
        align=center,
        thick
    },
    arr/.style={
        -{Stealth[length=4pt]},
        thick,
        tscgblue!70
    },
    phase/.style={
        font=\tiny\sffamily\color{tscggray},
        above=1pt
    }
]

% Input
\node[iobox] (input) {Raw\\Prompt};

% Phase 1: Parse
\node[box, right=of input] (parse) {1. Parse};

% Phase 2: Compression
\node[box, right=of parse] (sdm) {2. SDM};
\node[box, right=of sdm] (tas) {3. TAS};
\node[box, right=of tas] (dro) {4. DRO};

% Phase 3: Structural
\node[box, right=of dro] (cfl) {5. CFL};
\node[box, right=of cfl] (cfo) {6. CFO};

% Phase 4: Fragility
\node[box, right=of cfo] (cas) {7. CAS};
\node[box, right=of cas] (sadf) {8. SAD-F};

% Phase 5: Closure + Emit
\node[box, right=of sadf] (ccp) {9. CCP};
\node[iobox, right=of ccp] (emit) {10. Emit};

% Output
\node[iobox, right=of emit] (output) {TSCG\\Output};

% Arrows
\draw[arr] (input) -- (parse);
\draw[arr] (parse) -- (sdm);
\draw[arr] (sdm) -- (tas);
\draw[arr] (tas) -- (dro);
\draw[arr] (dro) -- (cfl);
\draw[arr] (cfl) -- (cfo);
\draw[arr] (cfo) -- (cas);
\draw[arr] (cas) -- (sadf);
\draw[arr] (sadf) -- (ccp);
\draw[arr] (ccp) -- (emit);
\draw[arr] (emit) -- (output);

% Phase labels
\node[phase] at (parse.north) {Parse};
\node[phase] at ($(sdm.north)!0.5!(dro.north)$) {Compression};
\node[phase] at ($(cfl.north)!0.5!(cfo.north)$) {Structural};
\node[phase] at ($(cas.north)!0.5!(sadf.north)$) {Fragility};
\node[phase] at ($(ccp.north)!0.5!(emit.north)$) {Closure};

% Phase grouping brackets
\draw[decorate, decoration={brace, amplitude=4pt, mirror}, tscggray!50, thick]
    (sdm.south west) ++(0,-0.15) -- (dro.south east |- sdm.south west) ++(0,-0.15)
    node[midway, below=4pt, font=\tiny\sffamily\color{tscggray}] {Density};

\draw[decorate, decoration={brace, amplitude=4pt, mirror}, tscggray!50, thick]
    (cfl.south west) ++(0,-0.15) -- (cfo.south east |- cfl.south west) ++(0,-0.15)
    node[midway, below=4pt, font=\tiny\sffamily\color{tscggray}] {Ordering};

\draw[decorate, decoration={brace, amplitude=4pt, mirror}, tscggray!50, thick]
    (cas.south west) ++(0,-0.15) -- (sadf.south east |- cas.south west) ++(0,-0.15)
    node[midway, below=4pt, font=\tiny\sffamily\color{tscggray}] {Anchoring};

\end{tikzpicture}
\caption{The TSCG 10-pass transform pipeline. Transforms execute in fixed order from left to right. The pipeline is organized into five phases: \emph{Parse} segments the input into semantic units; \emph{Compression} (SDM, TAS, DRO) removes filler and optimizes delimiters; \emph{Structural} (CFL, CFO) positions constraints and orders operations; \emph{Fragility} (CAS, SAD-F) scores and reinforces critical parameters; \emph{Closure} (CCP, Emit) appends the semantic closure block and generates the final output. Each transform is a pure function; the pipeline composition $\Pi = \tau_{10} \circ \cdots \circ \tau_1$ is deterministic.}
\label{fig:pipeline}
\end{figure*}
